% Bagian: incompleteness
% Bab: theories-computability
% Subbagian: introduction

\documentclass[../../../include/open-logic-section]{subfiles}

\begin{document}

\olfileid[id]{inc}{tcp}{int}

\olsection{Pengantar}

\begin{editorial}
Bagian ini perlu ditulis ulang.
\end{editorial}

Kita telah memperoleh hal-hal berikut:
\begin{enumerate}
\item definisi tentang apa artinya suatu fungsi dapat direpresentasikan
  dalam~$\Th{Q}$ (\olref[req][int]{defn:representable-fn});
\item definisi tentang apa artinya suatu relasi dapat direpresentasikan
  dalam~$\Th{Q}$ (\olref[req][rel]{defn:representing-relations});
\item teorema yang menyatakan bahwa fungsi-fungsi yang dapat direpresentasikan
  dalam~$\Th{Q}$ tepat merupakan fungsi-fungsi yang dapat dihitung
  (\olref[req][int]{thm:representable-iff-comp});
\item teorema yang menyatakan bahwa relasi-relasi yang dapat direpresentasikan
  dalam~$\Th{Q}$ tepat merupakan relasi-relasi yang dapat dihitung
  \olref[req][rel]{thm:representing-rels}.
\end{enumerate}

Suatu {\em teori} adalah himpunan kalimat yang tertutup secara deduktif,
yakni memiliki sifat bahwa setiap kali $T$ membuktikan~$!A$, maka $!A$ berada
dalam~$T$. Barangkali paling baik memandang teori sebagai kumpulan kalimat
beserta semua hal yang diimplikasikan oleh kalimat-kalimat tersebut. Mulai
sekarang, kita akan menggunakan~$\Th{Q}$ untuk menyebut {\em teori} yang
terdiri atas himpunan kalimat yang dapat diturunkan dari delapan aksioma dalam
\olref[req][int]{sec}. Ingat bahwa kita dapat mengodekan formula-formula
$\Th{Q}$ sebagai bilangan; jika $!A$ adalah formula semacam itu, misalkan
$\Gn{!A}$ menyatakan bilangan yang mengodekan~$!A$. Dengan pengodean ini, kita
sekarang dapat menanyakan apakah berbagai himpunan formula dapat dihitung atau
tidak.

\end{document}
